\documentclass[12pt,a4paper]{article}
\usepackage[utf8]{inputenc}
\usepackage[T1]{fontenc}
\usepackage[turkish]{babel}
\usepackage{geometry}
\usepackage{booktabs}
\usepackage{longtable}
\usepackage{array}
\usepackage{xcolor}
\usepackage{listings}
\usepackage{hyperref}
\usepackage{titlesec}
\usepackage{enumitem}
\usepackage{fancyhdr}
\usepackage{graphicx}

\geometry{margin=2.5cm}

\definecolor{primary}{RGB}{224,122,53}
\definecolor{codeblue}{RGB}{0,80,160}
\definecolor{codegray}{RGB}{100,100,100}
\definecolor{codebg}{RGB}{248,248,248}
\definecolor{successgreen}{RGB}{34,139,34}
\definecolor{errorred}{RGB}{180,0,0}

\lstset{
  backgroundcolor=\color{codebg},
  basicstyle=\ttfamily\small,
  keywordstyle=\color{codeblue}\bfseries,
  commentstyle=\color{codegray}\itshape,
  stringstyle=\color{primary},
  breaklines=true,
  frame=single,
  rulecolor=\color{primary!40},
  xleftmargin=8pt,
  xrightmargin=8pt,
  aboveskip=6pt,
  belowskip=6pt,
}

\titleformat{\section}{\large\bfseries\color{primary}}{}{0em}{}[\titlerule]
\titleformat{\subsection}{\normalsize\bfseries\color{codeblue}}{}{0em}{}

\pagestyle{fancy}
\fancyhf{}
\rhead{\textcolor{primary}{\textbf{CanYoldaşı}}}
\lhead{RevenueCat Entegrasyon Raporu}
\rfoot{\thepage}
\lfoot{\textcolor{codegray}{\small 16 Temmuz 2026}}

\hypersetup{
  colorlinks=true,
  linkcolor=primary,
  urlcolor=codeblue,
}

% ────────────────────────────────────────────
\begin{document}

\begin{titlepage}
  \centering
  \vspace*{3cm}
  {\Huge\bfseries\color{primary} CanYoldaşı\par}
  \vspace{0.5cm}
  {\Large RevenueCat Entegrasyon Raporu\par}
  \vspace{0.3cm}
  {\large Adım 4 · Adım 8 · Adım 9\par}
  \vspace{1.5cm}
  \hrule height 1pt
  \vspace{0.4cm}
  {\normalsize
    \textbf{Proje:} CanYoldaşı — Sokak Hayvanı Mobil Uygulaması\\[4pt]
    \textbf{Paket:} \texttt{react-native-purchases@10.4.2}\\[4pt]
    \textbf{Expo SDK:} 54.0.27 \quad \textbf{React Native:} 0.81.5\\[4pt]
    \textbf{Tarih:} 16 Temmuz 2026
  }
  \vspace{0.4cm}
  \hrule height 1pt
  \vfill
  {\small\color{codegray}
    Bu belge CanYoldaşı mobil uygulamasına yapılan RevenueCat\\
    in-app purchase entegrasyonunu üç aşamalı olarak belgelemektedir.
  }
\end{titlepage}

% ────────────────────────────────────────────
\tableofcontents
\newpage

% ════════════════════════════════════════════
\section{Genel Bakış}

Bu rapor, CanYoldaşı mobil uygulamasına RevenueCat Test Store entegrasyonunu
kapsayan üç uygulama adımını belgelemektedir.
Her adım bağımsız olarak planlanmış, uygulanmış ve TypeScript denetimiyle doğrulanmıştır.

\begin{center}
\begin{tabular}{lll}
  \toprule
  \textbf{Adım} & \textbf{Konu} & \textbf{Durum}\\
  \midrule
  4  & SDK Başlatma ve Kullanıcı Kimliği & \textcolor{successgreen}{\bfseries PASS}\\
  8  & Offerings Yükleme ve Paket Gösterimi & \textcolor{successgreen}{\bfseries PASS}\\
  9  & Test Store Satın Alma Akışı & \textcolor{successgreen}{\bfseries PASS}\\
  \bottomrule
\end{tabular}
\end{center}

\medskip
\noindent\textbf{TypeScript sonucu tüm adımlarda:} 0 hata, \texttt{any} / \texttt{@ts-ignore} / \texttt{@ts-expect-error} yok.

% ════════════════════════════════════════════
\section{Adım 4 — SDK Başlatma ve Kullanıcı Kimliği}

\subsection{Hedef}

RevenueCat SDK'sını uygulamanın tek bir merkezi noktasından idempotent biçimde başlatmak;
kullanıcı kimliğini JWT auth lifecycle'ına bağlamak.

\subsection{Oluşturulan ve değiştirilen dosyalar}

\begin{itemize}[leftmargin=1.5em]
  \item \texttt{artifacts/mobile/services/revenueCat.ts} — \textbf{yeni dosya}
  \item \texttt{artifacts/mobile/contexts/AuthContext.tsx}
  \item \texttt{artifacts/mobile/app/\_layout.tsx}
\end{itemize}

\subsection{Merkezi Servis: \texttt{revenueCat.ts}}

\begin{lstlisting}[language={}]
// Idempotency guard — Purchases.configure() en fazla bir kez çağrılır
let _initialized = false;
let _currentRcUserId: string | null = null;

export async function initializeRevenueCat(userId?: string | null): Promise<void>
export async function loginRevenueCat(userId: string): Promise<void>
export async function logoutRevenueCat(): Promise<void>
export function isRevenueCatInitialized(): boolean
\end{lstlisting}

\subsubsection{İdempotency Garantisi}

\texttt{Purchases.configure()} yalnızca \texttt{\_initialized === false} iken çağrılır.
React Strict Mode, Expo Fast Refresh ve bileşen remount durumlarında güvenlidir.

\subsubsection{Ortam Değişkeni}

\begin{lstlisting}[language={}]
const apiKey = process.env.EXPO_PUBLIC_REVENUECAT_API_KEY;
if (!apiKey) {
  console.warn("[RevenueCat] API key is missing");
  return; // configure() çağrılmaz
}
\end{lstlisting}

API anahtarı kaynak koda gömülmez; eksikse uygulama çökmez, sadece uyarı verilir.

\subsubsection{Debug Loglama}

\begin{lstlisting}[language={}]
if (__DEV__) {
  Purchases.setLogLevel(LOG_LEVEL.DEBUG); // configure()'dan ÖNCE
}
\end{lstlisting}

\subsubsection{Web Platform Koruması}

Her fonksiyonun başında \texttt{Platform.OS === "web"} kontrolü yapılır ve erken dönülür.
SDK \texttt{await import("react-native-purchases")} ile lazy yüklenir, web build çökmez.

\subsection{AuthContext Entegrasyonu}

\begin{center}
\begin{tabular}{ll}
  \toprule
  \textbf{Auth Olayı} & \textbf{RevenueCat Çağrısı}\\
  \midrule
  Uygulama başlangıcı — oturum var & \texttt{initializeRevenueCat(user.id)}\\
  Uygulama başlangıcı — oturum yok & \texttt{initializeRevenueCat(null)} (anonim)\\
  Kayıt ol (\texttt{register}) & \texttt{loginRevenueCat(user.id)}\\
  Giriş yap (\texttt{login}) & \texttt{loginRevenueCat(user.id)}\\
  Çıkış yap (\texttt{logout}) & \texttt{logoutRevenueCat()} → anonim sıfırlama\\
  \bottomrule
\end{tabular}
\end{center}

Tüm çağrılar \texttt{.catch(() => \{\})} ile non-blocking; RC hatası auth akışını bozmaz.

\subsection{\_layout.tsx Temizliği}

Eski IIFE başlatma bloğu kaldırıldı:
\begin{lstlisting}[language={}]
// KALDIRILDI — artık services/revenueCat.ts tek kaynak
if (Platform.OS !== "web") {
  (async () => { ... Purchases.configure(...) ... })();
}
\end{lstlisting}

\subsection{Uyarı: Test Store Anahtarı}

\begin{lstlisting}[language={}]
// ⚠️ EXPO_PUBLIC_REVENUECAT_API_KEY şu an TEST STORE anahtarıdır.
// App Store / Google Play üretimine geçmeden önce
// platform-specific üretim anahtarlarıyla değiştirilmelidir.
// Test Store anahtarıyla üretim yapılmaz.
\end{lstlisting}

% ════════════════════════════════════════════
\section{Adım 8 — RevenueCat Offerings Yükleme}

\subsection{Hedef}

\texttt{canyoldasi\_boost} offering'ini RC'den yüklemek;
üç paketi lokalize başlık ve fiyatla göstermek;
yükleme, hata ve boş durum UI'larını uygulamak.

\subsection{Değiştirilen dosyalar}

\begin{itemize}[leftmargin=1.5em]
  \item \texttt{artifacts/mobile/services/revenueCat.ts} — \texttt{fetchOfferings()} eklendi
  \item \texttt{artifacts/mobile/app/boost-packages.tsx} — RC offerings entegrasyonu
\end{itemize}

\subsection{fetchOfferings() Fonksiyonu}

\begin{lstlisting}[language={}]
export async function fetchOfferings(): Promise<PurchasesPackage[]> {
  // Önce current offering, sonra canyoldasi_boost fallback
  const offering = allOfferings.current
    ?? allOfferings.all["canyoldasi_boost"]
    ?? null;

  // Beklenmedik identifier uyarısı
  if (__DEV__ && offering.identifier !== "canyoldasi_boost") {
    console.warn(`[RevenueCat] Unexpected offering: "${offering.identifier}"`);
  }

  return offering.availablePackages as PurchasesPackage[];
}
\end{lstlisting}

\subsection{BOOST\_PACKAGE\_CONFIG — Merkezi Yapılandırma}

\begin{lstlisting}[language={}]
const BOOST_PACKAGE_CONFIG: Record<string, BoostPackageMeta> = {
  boost_1_day:  { durationDays: 1, isPopular: false, sortOrder: 0 },
  boost_3_days: { durationDays: 3, isPopular: true,  sortOrder: 1 },
  boost_7_days: { durationDays: 7, isPopular: false, sortOrder: 2 },
};
\end{lstlisting}

Magic string'ler tek kaynakta toplanır.
Sıralama \texttt{sortOrder}'a göre — RC'nin döndürdüğü sıra önemsizdir.
Bilinmeyen identifier'lar uygulamayı çökertmez.

\subsection{Gösterilen Paket Bilgileri}

\begin{center}
\begin{tabular}{lll}
  \toprule
  \textbf{Identifier} & \textbf{Süre} & \textbf{Fiyat Kaynağı}\\
  \midrule
  \texttt{boost\_1\_day}  & 1 gün & \texttt{pkg.product.priceString}\\
  \texttt{boost\_3\_days} & 3 gün (Popüler) & \texttt{pkg.product.priceString}\\
  \texttt{boost\_7\_days} & 7 gün & \texttt{pkg.product.priceString}\\
  \bottomrule
\end{tabular}
\end{center}

Fiyat kesinlikle hardcode değildir. Başlık \texttt{pkg.product.title} (RC lokalize).

\subsection{Durum Yönetimi}

\begin{center}
\begin{tabular}{ll}
  \toprule
  \textbf{Durum} & \textbf{UI}\\
  \midrule
  Yükleniyor & \texttt{ActivityIndicator} + "Paketler yükleniyor\ldots"\\
  Hata & "Öne çıkarma paketleri yüklenemedi." + Tekrar Dene\\
  Boş liste & "Şu an aktif öne çıkarma paketi bulunmuyor."\\
  Web & Uyarı banner + buton pasif\\
  \bottomrule
\end{tabular}
\end{center}

Duplike fetch koruması: \texttt{useRef(false)} ile React Strict Mode'da çift istek engellenir.

\subsection{Seçili Paket}

\begin{lstlisting}[language={}]
const [selectedPackage, setSelectedPackage] =
  useState<PurchasesPackage | null>(null);

// Varsayılan: boost_3_days; yoksa ilk paket; hiç yoksa null
const preferred =
  sorted.find(p => p.identifier === "boost_3_days")
  ?? sorted[0]
  ?? null;
\end{lstlisting}

% ════════════════════════════════════════════
\section{Adım 9 — Test Store Satın Alma Akışı}

\subsection{Hedef}

RevenueCat Test Store üzerinden gerçek satın alma akışını başlatmak;
durum makinesini uygulamak; DB aktivasyonunu henüz yapmamak.

\subsection{Değiştirilen dosyalar}

\begin{itemize}[leftmargin=1.5em]
  \item \texttt{artifacts/mobile/app/boost-packages.tsx}
\end{itemize}

\subsection{Purchase State Machine}

\begin{lstlisting}[language={}]
type PurchaseState =
  | "idle"       // Başlangıç / iptal / hata sonrası
  | "starting"   // Buton tıklandı, kontroller geçti
  | "processing" // purchasePackage() çalışıyor
  | "success"    // RC başarı döndürdü
  | "cancelled"  // Kullanıcı iptal etti
  | "pending"    // Mağaza işliyor
  | "error";     // RC hatası
\end{lstlisting}

\subsection{CompletedBoostPurchase — Tip Güvenli Sonuç Modeli}

\begin{lstlisting}[language={}]
type CompletedBoostPurchase = {
  packageIdentifier: string;
  productIdentifier: string;
  customerInfo: CustomerInfo; // react-native-purchases tipi
};
\end{lstlisting}

\subsection{Tip Güvenli Purchase Executor}

\begin{lstlisting}[language={}]
async function executePurchase(
  pkg: PurchasesPackage
): Promise<MakePurchaseResult> {
  const { default: Purchases } = await import("react-native-purchases");
  return Purchases.purchasePackage(pkg);
}
\end{lstlisting}

\texttt{any}, \texttt{@ts-ignore}, \texttt{@ts-expect-error} kullanılmamıştır.

\subsection{handleVerifiedPurchaseResult — Adım 10 Placeholder}

\begin{lstlisting}[language={}]
function handleVerifiedPurchaseResult(
  result: CompletedBoostPurchase
): void {
  if (__DEV__) {
    console.log("[Boost] Test Store purchase completed.",
      result.packageIdentifier, result.productIdentifier);
  }
  // TODO Adım 10: POST /api/boost/verify-iap → listing promosyon aktivasyonu
}
\end{lstlisting}

\subsection{Hata Sınıflandırma}

\begin{center}
\begin{tabular}{lll}
  \toprule
  \textbf{RC Hata Durumu} & \textbf{Tespit} & \textbf{Türkçe Mesaj}\\
  \midrule
  Kullanıcı iptali & \texttt{userCancelled === true} veya \texttt{code === 1} & "İptal edildi."\\
  Beklemede & \texttt{code === 6} & "Mağaza işliyor."\\
  Ürün yok & \texttt{code === 3} & "Satın alınamıyor."\\
  Ağ hatası & \texttt{code === 23} & "Bağlantı sorunu."\\
  Expo Go & Hata mesajı analizi & "Development Build gerekli."\\
  Diğer & — & "Lütfen tekrar deneyin."\\
  \bottomrule
\end{tabular}
\end{center}

Stack trace, receipt, token veya API key kullanıcıya gösterilmez.

\subsection{Güvenlik Kontrolleri}

\begin{center}
\begin{tabular}{ll}
  \toprule
  \textbf{Kontrol} & \textbf{Sonuç}\\
  \midrule
  \texttt{Platform.OS === "web"} & "iOS/Android uygulaması gerekli."\\
  \texttt{user === null} & "Giriş yapmanız gerekiyor."\\
  \texttt{selectedPackage === null} & "Paket seçin."\\
  \texttt{purchasingRef.current === true} & İkinci satın alma başlatılmaz\\
  \bottomrule
\end{tabular}
\end{center}

\subsection{Başarı UX (Bu Aşama İçin)}

\begin{lstlisting}[language={}]
Alert.alert(
  "Test Satın Alma Tamamlandı",
  "Test satın alma işlemi RevenueCat tarafından tamamlandı.\n\n"
  + "İlan aktivasyonu bir sonraki aşamada bağlanacak."
);
\end{lstlisting}

"İlan öne çıkarıldı" mesajı gösterilmez çünkü DB aktivasyonu henüz yapılmamıştır.

\subsection{Buton Durumu}

\begin{lstlisting}[language={}]
const isPurchasing = purchaseState === "starting"
                  || purchaseState === "processing";
const ctaDisabled  = isPurchasing || offeringsLoading
                  || !!offeringsError || !selectedPackage;
\end{lstlisting}

\begin{center}
\begin{tabular}{ll}
  \toprule
  \textbf{Durum} & \textbf{Buton Etiketi}\\
  \midrule
  Normal & "Öne Çıkarmayı Satın Al"\\
  Başlatılıyor & "Başlatılıyor\ldots"\\
  İşleniyor & "İşleniyor\ldots" + ActivityIndicator\\
  Web & "Uygulama Gerekli" (pasif)\\
  \bottomrule
\end{tabular}
\end{center}

% ════════════════════════════════════════════
\section{Test Senaryoları}

\begin{longtable}{p{4cm} p{5.5cm} l}
  \toprule
  \textbf{Senaryo} & \textbf{Beklenen Davranış} & \textbf{DB}\\
  \midrule
  \endhead
  Paket seç, satın al & RC Test Store akışı başlar & Değişmez\\
  Satın almayı iptal et & Yükleme sıfırlanır, ekran açık kalır & Değişmez\\
  Test Store başarısı & Dev-aşaması mesajı gösterilir & Değişmez\\
  Butona çift tıkla & Yalnızca 1 satın alma akışı başlar & Değişmez\\
  selectedPackage null & Satın alma başlamaz & Değişmez\\
  Giriş yapılmamış kullanıcı & Satın alma başlamaz & Değişmez\\
  Ürün erişilemiyor & Türkçe hata mesajı & Değişmez\\
  Ağ hatası & Türkçe hata, yükleme sıfırlanır & Değişmez\\
  Expo Go & Fake başarı yok, Build uyarısı & Değişmez\\
  Web & Native akış başlamaz, çökme yok & Değişmez\\
  \bottomrule
\end{longtable}

% ════════════════════════════════════════════
\section{Genel Durum}

\begin{center}
\begin{tabular}{ll}
  \toprule
  \textbf{Kontrol} & \textbf{Sonuç}\\
  \midrule
  TypeScript (tüm adımlar) & \textcolor{successgreen}{\bfseries PASS — 0 hata}\\
  \texttt{Purchases.configure} tekli kaynak & \textcolor{successgreen}{\bfseries Evet}\ (\texttt{services/revenueCat.ts})\\
  \texttt{any} kullanımı & \textcolor{successgreen}{\bfseries Yok}\\
  \texttt{@ts-ignore} / \texttt{@ts-expect-error} & \textcolor{successgreen}{\bfseries Yok}\\
  Hardcode fiyat & \textcolor{successgreen}{\bfseries Yok}\\
  Hardcode API key & \textcolor{successgreen}{\bfseries Yok}\\
  DB aktivasyonu (\texttt{promoted\_until}) & \textcolor{successgreen}{\bfseries Yok}\\
  Web build çökmesi & \textcolor{successgreen}{\bfseries Yok}\\
  \bottomrule
\end{tabular}
\end{center}

\subsection{Sonraki Adım (Adım 10)}

\texttt{handleVerifiedPurchaseResult()} fonksiyonu içinde \texttt{/api/boost/verify-iap}
endpoint'ine POST isteği gönderilecek; \texttt{listing\_promotions} tablosu güncellenecek;
ilan promosyon aktivasyonu gerçekleştirilecek.

\end{document}
